\documentclass[11pt]{article}

\usepackage[margin=1in]{geometry}
\usepackage{amsmath}
\usepackage{amssymb}
\usepackage{booktabs}
\usepackage{array}
\usepackage{tabularx}
\usepackage{xurl}
\usepackage{hyperref}
\usepackage{authblk}
\newcolumntype{P}[1]{>{\raggedright\arraybackslash}p{#1}}
\newcolumntype{Y}{>{\raggedright\arraybackslash}X}

\title{Five-Scale Audit of an e-5-137 Finite-Field Parameterization}

\author{I. Mikhailov}
\affil{Independent Researcher, Moscow, Russia}

\date{May 30, 2026}

\begin{document}

\maketitle

\begin{abstract}
This note records a finite-field parameterization built from $e$, $\pi$,
$\alpha^{-1}$, $N=5$, $D=26$, and arithmetic over GF(137).  A formula atlas
first lists standard reference equations from mechanics, electromagnetism,
thermodynamics, relativity, quantum mechanics, nuclear physics, and cosmology.
The project-specific material is then organized as a five-scale audit: MeV
nuclear scales, eV lepton and QED scales, molecular and semantic repair
scales, macroscopic gas and thermodynamic coarse-graining, and cosmological
closure tests.  The construction is phenomenological.  It is not a replacement
for the Standard Model, QCD, statistical mechanics, or $\Lambda$CDM.  Each
section therefore reports either a numerical benchmark or a null result.
\end{abstract}

\section{Scope}

The project uses a small set of integer and transcendental constants to define
testable formulae:
\begin{equation}
    N=5,\qquad D=26,\qquad p=137.
\end{equation}
The finite-field arithmetic uses residues in
\begin{equation}
    \mathbb{F}_{137}\simeq\mathbb{Z}/137\mathbb{Z}.
\end{equation}
The main dimensionless compression operator is
\begin{equation}
    q=\frac{N(N-2)}{e^4\pi^3}
      =0.008860611874290873.
\end{equation}
The phase difference used in several audits is
\begin{equation}
    \delta\phi=\cos(1/N)-\cos(2/N)
              =0.05900558383835652.
\end{equation}
The finite-field threshold used by the integer kernels and repair channels is
\begin{equation}
    I_5=2(D-N)=42.
\end{equation}

The purpose of the note is not to infer physics from numerology.  The purpose
is to record which formulae survive fixed tests and which do not.  A retained
formula must specify its domain, inputs, residuals, and failure modes.

This paper belongs to a companion series of e-5-137 and GF(137) audit
repositories.  The repositories intentionally share notation, implementation
patterns, and reproducibility conventions.  The scope of this archive is the
finite-field audit, the five-test matrix, and the English/Russian preprint
artifacts specific to this paper; related repositories should be cited
separately when their own results are used.

\section{Common Operators and Data Structures}

The computational layer uses three reusable objects.

First, a byte-valued residue vector stores model parameters:
\begin{equation}
    x,w,b\in\mathbb{F}_{137}.
\end{equation}
The edge-inference kernel computes
\begin{equation}
    R=(XW+b)\bmod 137,\qquad A_{ij}=\mathbf{1}_{R_{ij}\ge42}.
\end{equation}
This is an integer inference primitive, not a floating-point neural-network
approximation theorem.

Second, the fault-tolerance layer replicates a symbol $s$ over $26$ axes:
\begin{equation}
    r_a=s+m_a(k)\pmod{137},\qquad a=0,\ldots,25,
\end{equation}
where $m_a(k)$ is a deterministic axis mask.  Repair subtracts the mask and
selects the most frequent residue:
\begin{equation}
    \hat{s}
    =
    \operatorname*{arg\,max}_{v\in\mathbb{F}_{137}}
    \#\{a:r_a-m_a(k)\equiv v\pmod{137}\}.
\end{equation}
With at most $10$ corrupted axes, at least $16$ correct axes remain:
\begin{equation}
    26-10=16.
\end{equation}
This gives a fixed repair budget for the toy channel.

Third, the long-context compression benchmark maps a $120{,}000$ token object
to
\begin{equation}
    V_{\rm ref}=117
\end{equation}
semantic vectors, or
\begin{equation}
    \frac{120{,}000}{117}=1025.6410256410256
    \quad {\rm tokens/vector}.
\end{equation}
This is a synthetic benchmark tied to the repository tests.  It is not a
general lossless-compression theorem.

\section{Formula Atlas Across Physical Scales}

This section lists the standard formulae used as reference structure from
school mechanics through undergraduate field theory and cosmology.  The
finite-field parameterization is not used to replace these equations.  It is
used only as an audit layer when a discrete residue model, compression channel,
or numerical benchmark is explicitly defined.

\subsection{Dimensional and Conservation Prefilter}

Every formula in the audit must first pass dimensional analysis:
\begin{equation}
    [v]=LT^{-1},\qquad
    [a]=LT^{-2},\qquad
    [F]=MLT^{-2},\qquad
    [E]=ML^2T^{-2}.
\end{equation}
The conserved quantities used as organizing checks are
\begin{equation}
    \frac{dE}{dt}=0,\qquad
    \frac{dp}{dt}=0,\qquad
    \frac{dL}{dt}=0,
\end{equation}
when the corresponding time-translation, space-translation, and rotational
symmetries hold.  A finite-field expression that violates dimensions is not
retained as a physical formula.

\subsection{Kinematics and Newtonian Mechanics}

The kinematic definitions are
\begin{equation}
    v=\frac{dx}{dt},\qquad
    a=\frac{dv}{dt}=\frac{d^2x}{dt^2}.
\end{equation}
For constant acceleration,
\begin{align}
    x(t)&=x_0+v_0t+\frac{1}{2}at^2,\\
    v(t)&=v_0+at,\\
    v^2&=v_0^2+2a(x-x_0).
\end{align}
Newtonian dynamics uses
\begin{equation}
    F=ma,\qquad p=mv,\qquad J=\int F\,dt=\Delta p.
\end{equation}
The work and energy formulae are
\begin{equation}
    W=\int F\,dx,\qquad
    K=\frac{1}{2}mv^2,\qquad
    U_g=mgh,\qquad
    U_s=\frac{1}{2}kx^2.
\end{equation}
The Lagrangian form is
\begin{equation}
    L=T-V,\qquad
    \frac{d}{dt}\frac{\partial L}{\partial \dot{q}}
    -\frac{\partial L}{\partial q}=0.
\end{equation}
The e-5-137 layer has no role in these identities except as a possible
integer representation of a numerical simulator.

\subsection{Gravitation and Orbital Dynamics}

Newtonian gravity is
\begin{equation}
    F=G\frac{m_1m_2}{r^2},\qquad
    g=\frac{GM}{r^2},\qquad
    \Phi=-\frac{GM}{r}.
\end{equation}
For a circular orbit and escape from radius $r$,
\begin{equation}
    v_{\rm orb}=\sqrt{\frac{GM}{r}},
    \qquad
    v_{\rm esc}=\sqrt{\frac{2GM}{r}}.
\end{equation}
Kepler's third law is
\begin{equation}
    T^2=\frac{4\pi^2}{GM}a^3.
\end{equation}
The relativistic extension used later in the cosmology section is not derived
from this Newtonian block; it is a separate metric theory.

\subsection{Oscillations, Waves, and Optics}

The harmonic oscillator obeys
\begin{equation}
    m\ddot{x}+kx=0,\qquad
    \omega=\sqrt{\frac{k}{m}},\qquad
    T=\frac{2\pi}{\omega}.
\end{equation}
The one-dimensional wave equation is
\begin{equation}
    \frac{\partial^2 y}{\partial t^2}
    =
    v^2\frac{\partial^2 y}{\partial x^2},
    \qquad
    v=\lambda f.
\end{equation}
For geometric and wave optics,
\begin{equation}
    n=\frac{c}{v},\qquad
    n_1\sin\theta_1=n_2\sin\theta_2,
\end{equation}
\begin{equation}
    \frac{1}{f}=\frac{1}{d_o}+\frac{1}{d_i},
    \qquad
    m=-\frac{d_i}{d_o},
\end{equation}
and diffraction maxima satisfy
\begin{equation}
    d\sin\theta=m\lambda.
\end{equation}
These formulae provide the continuous wave reference for the discrete
phase-lane language used by the repository.

\subsection{Electricity and Magnetism}

The electrostatic force and field definitions are
\begin{equation}
    F=k_e\frac{q_1q_2}{r^2},\qquad
    E=\frac{F}{q},\qquad
    V=\frac{W}{q}.
\end{equation}
Circuit identities are
\begin{equation}
    I=\frac{dq}{dt},\qquad
    U=IR,\qquad
    P=UI=I^2R=\frac{U^2}{R}.
\end{equation}
For capacitors and inductors,
\begin{equation}
    C=\frac{Q}{U},\qquad
    E_C=\frac{1}{2}CU^2,\qquad
    E_L=\frac{1}{2}LI^2.
\end{equation}
The Lorentz force is
\begin{equation}
    F=q(E+v\times B).
\end{equation}
Maxwell's equations in SI units are
\begin{align}
    \nabla\cdot E&=\frac{\rho}{\varepsilon_0},&
    \nabla\cdot B&=0,\\
    \nabla\times E&=-\frac{\partial B}{\partial t},&
    \nabla\times B&=\mu_0J+\mu_0\varepsilon_0\frac{\partial E}{\partial t}.
\end{align}
The fine-structure constant enters the lepton parameterization through
$\alpha^{-1}$, but the Maxwell equations themselves are not modified.

\subsection{Thermodynamics and Statistical Mechanics}

The ideal-gas equation is
\begin{equation}
    PV=nRT=Nk_BT.
\end{equation}
The first and second laws are written as
\begin{equation}
    \Delta U=Q-W,\qquad
    dS\ge \frac{\delta Q}{T}.
\end{equation}
For heat capacity and phase transitions,
\begin{equation}
    Q=mc\Delta T,\qquad
    Q=\lambda m.
\end{equation}
The Carnot efficiency is
\begin{equation}
    \eta=1-\frac{T_c}{T_h}.
\end{equation}
The statistical form is
\begin{equation}
    S=k_B\ln\Omega,\qquad
    p_i=\frac{e^{-\beta E_i}}{Z},\qquad
    Z=\sum_i e^{-\beta E_i}.
\end{equation}
In this paper, GF(137) appears at this level only as a finite microscopic
alphabet whose randomized high-occupancy limit can be coarse-grained to the
usual Boltzmann distribution.

\subsection{Special Relativity}

The Lorentz factor is
\begin{equation}
    \gamma=\frac{1}{\sqrt{1-v^2/c^2}}.
\end{equation}
Time dilation, length contraction, and energy-momentum relations are
\begin{equation}
    \Delta t'=\gamma\Delta t,\qquad
    L'=\frac{L}{\gamma},
\end{equation}
\begin{equation}
    E=mc^2,\qquad
    E^2=p^2c^2+m^2c^4.
\end{equation}
The horizon-gated dark-energy audit uses a time-dilation-like factor, but its
numerical effect is recorded as a null result in the cosmology block.

\subsection{Quantum Mechanics}

The Planck and de Broglie relations are
\begin{equation}
    E=hf=\hbar\omega,\qquad
    p=\hbar k=\frac{h}{\lambda}.
\end{equation}
The uncertainty relation is
\begin{equation}
    \Delta x\,\Delta p\ge\frac{\hbar}{2}.
\end{equation}
The Schr\"odinger equation is
\begin{equation}
    i\hbar\frac{\partial\psi}{\partial t}=\hat{H}\psi,
    \qquad
    \hat{H}=-\frac{\hbar^2}{2m}\nabla^2+V.
\end{equation}
For hydrogenic levels,
\begin{equation}
    E_n=-13.6~{\rm eV}\frac{Z^2}{n^2}.
\end{equation}
The project uses $5e~{\rm eV}=13.591409142295225~{\rm eV}$ as a numerical
anchor near the hydrogen scale.  This is a parameterization choice, not a
derivation of QED.

\subsection{Nuclear and Particle Bookkeeping}

Nuclear binding and radioactive decay use
\begin{equation}
    E_b=\Delta mc^2,\qquad
    R=R_0A^{1/3},
\end{equation}
\begin{equation}
    N(t)=N_0e^{-\lambda t},\qquad
    T_{1/2}=\frac{\ln2}{\lambda}.
\end{equation}
Reaction energy is
\begin{equation}
    Q=(m_{\rm initial}-m_{\rm final})c^2.
\end{equation}
The repository's nuclear calculation is tested against $E_b=\Delta mc^2$ and
fails for helium, carbon, and iron at the $<10$ ppm criterion.

\subsection{Cosmological Observables}

The Hubble relation and redshift definition are
\begin{equation}
    v=H_0d,\qquad
    1+z=\frac{a_0}{a}.
\end{equation}
The critical density is
\begin{equation}
    \rho_c=\frac{3H^2}{8\pi G}.
\end{equation}
The Friedmann equation is
\begin{equation}
    H^2=\frac{8\pi G}{3}\rho-\frac{kc^2}{a^2}
    +\frac{\Lambda c^2}{3}.
\end{equation}
The equation-of-state parameter is
\begin{equation}
    w=\frac{p}{\rho c^2}.
\end{equation}
The effective radiation count is
\begin{equation}
    N_{\rm eff}=3.046+\Delta N_{\rm eff}.
\end{equation}
The model-specific CMB audit uses
\begin{equation}
    \Delta N_{\rm eff}
    =
    \frac{D}{137}\delta\phi(1-q)
    =
    0.011098917626279356.
\end{equation}

\subsection{Atlas Status Matrix}

\begin{table}[h]
\centering
\caption{Reference formula classes and the role of the finite-field layer.}
\scriptsize
\setlength{\tabcolsep}{3pt}
\begin{tabularx}{\linewidth}{@{}P{0.12\linewidth}P{0.17\linewidth}P{0.22\linewidth}P{0.20\linewidth}Y@{}}
\toprule
Domain & Core formula & Conservation or limit & GF(137) role & Status \\
\midrule
Mechanics & $F=ma$, $L=T-V$ & energy, momentum, angular momentum & numerical representation only & reference theory \\
Gravity & $F=Gm_1m_2/r^2$ & central-force orbits & none & reference theory \\
Waves/optics & $v=\lambda f$, Snell, lens law & phase propagation & phase-lane analogy & reference \\
EM & Maxwell equations, Lorentz force & charge and field continuity & $\alpha^{-1}$ in lepton ansatz only & reference \\
Thermal & $PV=nRT$, $S=k_B\ln\Omega$ & maximum entropy limit & finite alphabet regulator & coarse-grained bridge \\
Quantum & $i\hbar\partial_t\psi=\hat{H}\psi$ & unitary evolution & $5e$ eV anchor & parameterization \\
Nuclear & $E_b=\Delta mc^2$ & mass-energy accounting & tested mass-defect encoder & null result \\
Cosmology & Friedmann, $N_{\rm eff}$, $w$ & homogeneous expansion & CMB and DESI audits & conditional test/null \\
\bottomrule
\end{tabularx}
\end{table}

\section{Five-Scale Physical and Computational Audit}

\subsection{Level 1: MeV Nuclear and QCD Scale}

The MeV-scale audit tests whether the $D=26$, $42$, and $117$ operators
predict nuclear mass defects.  The tested form is
\begin{equation}
    \Delta M(Z,A)
    =
    Z^2E_{\rm step}(Z)
    \left[
        1+\frac{42}{137D}\left(1-\frac{13}{117}\right)
    \right],
\end{equation}
where the running step is
\begin{equation}
    E_{\rm step}(1)=5e~{\rm eV},
\end{equation}
and
\begin{equation}
    E_{\rm step}(Z>1)
    =
    5e\,(137)(42)
    \left(1-\frac{D}{137N}\right){\rm eV}
    =
    0.07523660444808945~{\rm MeV}.
\end{equation}
The fixed correction factor is
\begin{equation}
    1+\frac{42}{137D}\left(1-\frac{13}{117}\right)
    =
    1.010481003181733.
\end{equation}

Reference nuclear masses are obtained from neutral isotope masses after
subtracting electron masses using CODATA constants
\cite{nistatomicweights,codata2018}.  The audit table is:
\begin{table}[h]
\centering
\caption{Nuclear mass-defect audit.  The ppm residual is computed from the
predicted nuclear mass relative to the tabulated nuclear mass.}
\begin{tabular}{lrrrr}
\toprule
Isotope & $Z$ & Reference defect MeV & Model defect MeV & Residual ppm \\
\midrule
$^1{\rm H}$     & 1  & $0.0000133172$ & $0.0000137339$ & $-0.000444$ \\
$^4{\rm He}$    & 2  & $28.2956897122$ & $0.3041006382$ & $7509.723753$ \\
$^{12}{\rm C}$  & 6  & $92.1618167441$ & $2.7369057434$ & $8002.327108$ \\
$^{56}{\rm Fe}$ & 26 & $492.259569845$ & $51.3930078482$ & $8463.590834$ \\
\bottomrule
\end{tabular}
\end{table}

For iron,
\begin{equation}
    Z_{\rm Fe}=26=D,
\end{equation}
but the numerical mass-defect gate fails:
\begin{equation}
    \Delta M_{\rm Fe}^{\rm model}=51.393007848156245~{\rm MeV},
\end{equation}
\begin{equation}
    \Delta M_{\rm Fe}^{\rm ref}=492.2595698447403~{\rm MeV},
\end{equation}
\begin{equation}
    \epsilon_{\rm Fe}=8463.590833506674~{\rm ppm}.
\end{equation}
The $<10$ ppm criterion is not satisfied for helium, carbon, or iron.  The
MeV-scale result is therefore negative.  The finite-field operators do not
derive nuclear binding energies.  QCD confinement, shell structure, pairing,
and many-body nuclear dynamics remain outside this parameterization.

\subsection{Level 2: eV/QED Lepton Scale}

The eV-scale branch uses the electron mass as the normalization anchor:
\begin{equation}
    m_e=0.51099895069~{\rm MeV},
    \qquad
    \alpha^{-1}=137.035999084.
\end{equation}
The base bookkeeping step is
\begin{equation}
    5e~{\rm eV}=13.591409142295225~{\rm eV}.
\end{equation}
The lepton mass ratios are written as generation-scale ansatz formulae:
\begin{equation}
    \frac{m_\mu}{m_e}
    =
    \alpha^{-1}\left(\frac{3}{2}+q\right)
    =
    206.768221426689,
\end{equation}
and
\begin{equation}
    \frac{m_\tau}{m_e}
    =
    \alpha^{-1}\left(N^2+qf_\tau\right),
\end{equation}
with
\begin{equation}
    f_\tau=42+2e^{-2}+2e^{-7}.
\end{equation}
The resulting tau/electron ratio is
\begin{equation}
    3477.2282035579738.
\end{equation}

The fourth charged benchmark is
\begin{equation}
    \frac{m_{\tau'}}{m_e}
    =
    \alpha^{-1}\left[D^2+q(DN)\right]
    =
    92794.18434487357,
\end{equation}
or
\begin{equation}
    m_{\tau'}=47.417730830364825~{\rm GeV}.
\end{equation}
This is below current ordinary sequential charged-lepton limits
\cite{pdg2024}.  It is therefore not an allowed Standard-Model-like fourth
charged lepton.

The sterile neutral benchmark is
\begin{equation}
    m_{\nu_{\tau'}}
    =
    m_e
    \left(\frac{F_{26}}{\alpha^{-1}}\right)
    \left[
        1+\frac{DN}{\pi}\delta\phi(1-s)
    \right],
\end{equation}
where
\begin{equation}
    F_{26}=121393,\qquad s=\frac{\alpha}{2\pi}.
\end{equation}
The value is
\begin{equation}
    m_{\nu_{\tau'}}=1.5566463427697075~{\rm GeV}.
\end{equation}
This is a mass benchmark only.  Relic abundance, lifetime, production history,
mixing angles, and detection limits require a separate particle model
\cite{adhikari2017}.

\subsection{Level 3: Molecular, Repair, and Semantic Scale}

Molecular binding energies lie in the eV range, but the repository does not
contain an electronic-structure solver.  The finite-field contribution at this
scale is a symbolic repair and compression channel rather than a molecular
Hamiltonian.

The threshold $42$ is used as a discrete activation and repair marker:
\begin{equation}
    A=\mathbf{1}_{R\ge42}.
\end{equation}
For the DNA toy channel, a symbol is encoded into $26$ masked residues.  If at
most $10$ axes are corrupted, the majority repair rule recovers the symbol in
the tested model:
\begin{equation}
    d_{\rm damaged}\le10
    \quad\Rightarrow\quad
    26-d_{\rm damaged}\ge16.
\end{equation}
The test suite verifies this condition for synthetic DNA replication and for a
semantic archive.  In the semantic compression benchmark,
\begin{equation}
    120{,}000~{\rm tokens}
    \mapsto
    117~{\rm vectors}
    \mapsto
    117\times130~{\rm bytes}.
\end{equation}
The repair channel is exact under the configured damage budget.  This is a
statement about the test harness.  It is not a proof of biological aging,
organism-level longevity, or chemical stability.

\subsection{Level 4: Macro/Gas Thermodynamic Limit}

The macroscopic gas scale is treated as a coarse-graining limit of a finite
alphabet.  Let each microscopic degree of freedom carry a residue
$a\in\mathbb{F}_{137}$.  Map it to a centered integer
\begin{equation}
    \rho(a)\in\{-68,-67,\ldots,67,68\}.
\end{equation}
For a three-component velocity surrogate, use
\begin{equation}
    v=(\rho(a_1),\rho(a_2),\rho(a_3))v_0,
\end{equation}
and assign kinetic energy
\begin{equation}
    \varepsilon(v)=\frac{m}{2}|v|^2.
\end{equation}
For $M$ independent cells, the maximum-entropy distribution under fixed mean
energy is obtained by maximizing
\begin{equation}
    S=-k_B\sum_i p_i\ln p_i
\end{equation}
subject to
\begin{equation}
    \sum_i p_i=1,\qquad \sum_i p_i\varepsilon_i=\bar{E}.
\end{equation}
The Lagrange multiplier calculation gives
\begin{equation}
    p_i=\frac{e^{-\beta\varepsilon_i}}{Z_1(\beta)},
    \qquad
    Z_1(\beta)=\sum_i e^{-\beta\varepsilon_i}.
\end{equation}
Thus the discrete residue alphabet becomes a finite regularization of the
Boltzmann distribution.

In the isotropic high-occupancy limit, the residue sum is replaced by a
velocity integral.  For $N_g$ particles in volume $V$,
\begin{equation}
    Z_{N_g}
    =
    \frac{1}{N_g!}
    \left[
        \frac{V}{h^3}
        \int_{\mathbb{R}^3}e^{-\beta p^2/(2m)}\,d^3p
    \right]^{N_g}
    =
    \frac{1}{N_g!}
    \left[
        V\left(\frac{2\pi m}{\beta h^2}\right)^{3/2}
    \right]^{N_g}.
\end{equation}
The pressure follows from
\begin{equation}
    P=\frac{1}{\beta}\frac{\partial \ln Z_{N_g}}{\partial V}
      =
    \frac{N_g}{\beta V}.
\end{equation}
Since $\beta^{-1}=k_BT$,
\begin{equation}
    PV=N_gk_BT=nRT.
\end{equation}
The Boltzmann entropy is the corresponding macrostate count:
\begin{equation}
    S=k_B\ln\Omega.
\end{equation}

The result is a standard statistical-mechanics derivation written with a
finite residue alphabet as the microscopic regulator.  The ideal-gas law is
not unique to GF(137).  The role of GF(137) is to give a bounded discrete
state space.  When phase information is randomized and the distribution becomes
isotropic, the modular order is washed out and the continuum thermodynamic
limit recovers the Mendeleev--Clapeyron equation and Boltzmann entropy
\cite{boltzmann1896,gibbs1902}.

\subsection{Level 5: Cosmological Scale}

The cosmological branch contains three tests.

First, the Hubble reference values used by the repository are
\begin{equation}
    H_{0,\rm Planck}=67.4,
    \qquad
    H_{0,\rm SH0ES}=73.04,
\end{equation}
with fractional difference
\begin{equation}
    \frac{73.04-67.4}{67.4}=0.0836795252225519.
\end{equation}
The same section records a phenomenological CPL pair
\begin{equation}
    w_0=-0.7519948527423932,
    \qquad
    w_a=-0.8600649695978589.
\end{equation}

Second, the Lagrangian dark-energy audit tested the two-field background
\begin{equation}
    \mathcal{L}_{\rm dark}
    =
    -\frac{1}{2}(\partial\phi)^2
    +\frac{1}{2}(\partial\chi)^2
    -\Lambda^4U(x,y).
\end{equation}
The horizon time-dilation gate was
\begin{equation}
    H_{\rm gate}(z,x)
    =
    \left[
        1-\frac{q}{117}e^z\cos^2(5x)
    \right]^{-1/2}.
\end{equation}
Its effect was too small:
\begin{equation}
    \frac{q}{117}=7.57317\times10^{-5},
    \qquad
    \max H_{\rm gate}\simeq1.00028.
\end{equation}
The accepted best row after the gate was
\begin{equation}
    w_0=-0.920331,\qquad w(z=1)=-1.00259.
\end{equation}
The target value was
\begin{equation}
    w_{\rm target}(z=1)=-1.18202733754,
\end{equation}
leaving
\begin{equation}
    w(z=1)-w_{\rm target}(z=1)=0.179442.
\end{equation}
The DESI-scale field-dynamics branch is therefore retained as an active
null result, not as a fitted phantom-dark-energy model.

Third, the CMB closure invariant is
\begin{equation}
    \Delta N_{\rm eff}
    =
    \frac{D}{137}\delta\phi(1-q)
    =
    0.011098917626279356,
\end{equation}
so that
\begin{equation}
    N_{\rm eff}=3.046+\Delta N_{\rm eff}
    =
    3.0570989176262793.
\end{equation}
Relative to the Planck 2018 CMB+BAO value
$2.99\pm0.17$ \cite{planck2018}, this is
\begin{equation}
    0.39469951544870036
\end{equation}
standard deviations from the central value.  This is consistent with the
absence of a large extra light free-streaming component in current extended
CMB fits \cite{actdr6extended}.

For the sterile benchmark mass
\begin{equation}
    m_{\nu_{\tau'}}=1.5566463427697075~{\rm GeV},
\end{equation}
the thermal-equivalent cutoff estimate
\begin{equation}
    \lambda_{\rm cut}
    \simeq0.11\left(\frac{m}{\rm keV}\right)^{-4/3}{\rm Mpc}
\end{equation}
gives
\begin{equation}
    \lambda_{\rm cut}
    =
    6.097325934119188\times10^{-7}~{\rm kpc}.
\end{equation}
This is far below the $100~{\rm kpc}$ guardrail and far below DESI/BAO
scales.  The conclusion is conditional on a sterile, cold or sufficiently
nonthermal production history.

\section{Applied Integer Implementation}

The edge-AI branch uses the same residue arithmetic for inference.  The
measured model sizes were
\begin{equation}
    M_{\rm GF(137)}=1.62793~{\rm KB},
    \qquad
    M_{\rm float32}=6.50391~{\rm KB},
\end{equation}
or a factor of approximately $4$ in storage.  For $1000$ forward passes, the
native ARM NEON loop took $12.1055$ ms against $58.7784$ ms for the NumPy
float32 baseline:
\begin{equation}
    \frac{58.7784}{12.1055}\simeq4.86.
\end{equation}
The call-through C++ path took $40.1530$ ms, giving a speedup of about
$1.46\times$ while retaining byte-valued weights.

\begin{table}[h]
\centering
\caption{Edge-inference storage and runtime audit.}
\begin{tabular}{lrrl}
\toprule
Backend & Memory KB & Time for $1000$ passes & Role \\
\midrule
float32 NumPy MLP & 6.50391 & 58.7784 ms & host baseline \\
GF(137) C++ fused \texttt{uint8\_t} & 1.62793 & 40.1530 ms & call-through path \\
GF(137) C++ NEON native loop & 1.62793 & 12.1055 ms & native loop \\
\bottomrule
\end{tabular}
\end{table}

A hardware implementation would expose fixed operations
\begin{align}
    {\tt mac137}(a,b,c)&=(c+ab)\bmod137,\\
    {\tt inv137}(a)&=a^{-1}\bmod137,\quad a\ne0,\\
    {\tt ge42}(a)&=\mathbf{1}_{a\ge42}.
\end{align}
The relevant engineering property is bounded integer arithmetic.  The kernel
does not require a general divider in the inner inference loop; modular
reduction can be implemented with constants, subtractive correction, or a
small lookup/inverse block.

\section{Summary Table}

\begin{table}[h]
\centering
\caption{Five-scale status summary.}
\begin{tabular}{llll}
\toprule
Level & Scale & Main quantity & Status \\
\midrule
1 & MeV/QCD & Fe residual $8463.590833506674$ ppm & failed $<10$ ppm gate \\
2 & eV/QED & $m_{\nu_{\tau'}}=1.5566463427697075$ GeV & mass benchmark only \\
3 & molecular/semantic & $26$ axes, $10$-axis damage budget & toy repair passes \\
4 & gas/thermal & $PV=N_gk_BT=nRT$ & coarse-grained limit \\
5 & cosmology & $N_{\rm eff}=3.0570989176262793$ & within Planck window \\
\bottomrule
\end{tabular}
\end{table}

\section{Limitations}

The failed nuclear mass-defect gate is part of the result.  No correction term
is added to force agreement with NIST/CODATA isotope masses.  The lepton and
sterile-neutrino formulae are parameterizations rather than a Lagrangian
derivation.  The molecular and DNA sections describe a symbolic finite-field
repair code, not a biochemical model.  The gas section is standard statistical
mechanics with a finite residue alphabet used as a regulator.  The cosmology
section reports closure checks and null results under specified assumptions;
it does not fit DESI, Planck, ACT, or Fermi-LAT likelihoods.

\section{Conclusion}

The e-5-137 construction is useful in this repository as a compact
audit language.  It gives reproducible integer kernels, fixed repair
thresholds, several numerical mass benchmarks, and explicit negative results.
The five-scale organization separates successful engineering tests from
physics claims that remain unsupported.  The retained scientific standard is
the residual table, not the integer coincidence.

\begin{thebibliography}{11}

\bibitem{pdg2024}
S. Navas et al. (Particle Data Group),
``Review of Particle Physics,''
\emph{Phys. Rev. D} \textbf{110}, 030001 (2024), with 2025 update.

\bibitem{adhikari2017}
R. Adhikari et al.,
``A White Paper on keV Sterile Neutrino Dark Matter,''
\emph{JCAP} \textbf{01}, 025 (2017).

\bibitem{planck2018}
N. Aghanim et al. (Planck Collaboration),
``Planck 2018 results. VI. Cosmological parameters,''
\emph{Astron. Astrophys.} \textbf{641}, A6 (2020),
arXiv:1807.06209.

\bibitem{actdr6extended}
Atacama Cosmology Telescope Collaboration,
``The Atacama Cosmology Telescope: DR6 Constraints on Extended Cosmological
Models,''
arXiv:2503.14454.

\bibitem{nistatomicweights}
National Institute of Standards and Technology,
``Atomic Weights and Isotopic Compositions with Relative Atomic Masses,''
\url{https://www.nist.gov/pml/atomic-weights-and-isotopic-compositions}.

\bibitem{codata2018}
National Institute of Standards and Technology,
``CODATA Recommended Values of the Fundamental Physical Constants: 2018,''
\url{https://pml.nist.gov/cuu/pdf/wall_2018.pdf}.

\bibitem{kuratov2025}
Y. Kuratov, M. Arkhipov, A. Bulatov, and M. Burtsev,
``Cramming 1568 Tokens into a Single Vector and Back Again: Exploring the
Limits of Embedding Space Capacity,''
arXiv:2502.13063 (2025).

\bibitem{boltzmann1896}
L. Boltzmann,
\emph{Vorlesungen \"uber Gastheorie},
J. A. Barth, Leipzig (1896).

\bibitem{gibbs1902}
J. W. Gibbs,
\emph{Elementary Principles in Statistical Mechanics},
Charles Scribner's Sons, New York (1902).

\end{thebibliography}

\end{document}
